% Related Work section for EpiContext paper
\section{Related Work}

\EpiContext{} sits at the intersection of three active research areas: agent memory management, context window optimization, and biologically-inspired AI systems. We review each in turn and position our contributions relative to existing work.

\subsection{Agent Memory Management}

Memory is a foundational capability for LLM agents, enabling them to maintain coherence across extended interactions~\citep{wang2024survey, hu2025survey}. Early work established the basic memory taxonomy of short-term (in-context) and long-term (external storage) memory~\citep{park2023generative}. \ReAct{}~\citep{yao2023react} introduced the thought-action-observation loop that remains the dominant agent paradigm, where each turn's output becomes part of the next turn's context. \Reflexion{}~\citep{shinn2023reflexion} enhanced this with verbal self-reflection, storing evaluation feedback as additional memory.

For long-term memory, \MemGPT{}~\citep{packer2023memgpt} pioneered the operating system analogy, treating context windows as virtual memory with paging between fast (in-context) and slow (external) storage. MemoryBank~\citep{zhong2024memorybank} incorporated psychological principles like the Ebbinghaus forgetting curve to model memory decay. A-MEM~\citep{xu2025amem} adopted the Zettelkasten note-taking method, organizing memories as an interconnected knowledge network with LLM-driven linking. ProMem~\citep{chen2025promem} identified the ``missing information accumulation'' problem in summary-based methods and proposed proactive self-questioning for more complete memory extraction. LightMem~\citep{sinha2025lightmem} explored using small language models (SLMs) for lightweight online memory control, trading some accuracy for substantial efficiency gains.

A critical limitation of these approaches is their treatment of memory as a \textit{passive store}---information is written, compressed, and retrieved, but the system does not actively learn which memories are valuable for which contexts. \EpiContext{} addresses this gap through fitness-driven epigenetic tags that continuously adapt based on task outcomes, enabling the system to learn optimal memory activation patterns over time.

\subsection{Context Window Optimization}

As LLM context windows have expanded from 4K to 2M+ tokens, a parallel line of work has focused on making effective use of this capacity. Liu et al.~\citep{liu2024lost} demonstrated that LLM performance follows a U-shaped curve with respect to information position, with content in the middle of long contexts being systematically underutilized. This finding motivated research into context compression and selective attention mechanisms.

Anthropic's context engineering framework~\citep{anthropic2025context} introduced progressive disclosure and dynamic context discovery as design principles, advocating that agents should pull relevant context on demand rather than receiving everything upfront. Cursor's dynamic context discovery~\citep{cursor2025dynamic} applies this principle in practice, treating tool responses as files and loading only necessary MCP tools. JetBrains Research~\citep{jetbrains2025context} systematically studied context management strategies across different agent architectures, finding that LLM-based summarization provides the best cost-quality trade-off but noting that context management remains ``treated as an engineering detail rather than a core research problem.''

For tool-specific optimization, AutoTool~\citep{zhang2025autotool} introduced dynamic tool selection through a Tool Inertia Graph, transforming the combinatorial action space into a constrained graph search problem. The Model Context Protocol (MCP)~\citep{mcp2025} standardized tool interfaces but left tool selection strategy to individual implementations.

\EpiContext{} advances beyond these approaches in two key ways. First, it unifies memory and tool management under a single epigenetic framework, enabling coordinated optimization rather than independent tuning. Second, it introduces a principled fitness function that provides a clear optimization objective, moving context management from heuristic engineering to formal optimization.

\subsection{Biologically-Inspired AI}

Biological systems have long served as inspiration for AI architectures, from artificial neural networks themselves to more recent neuromorphic computing~\citep{schuman2017survey} and spiking neural networks~\citep{tavanaei2019deep}. In the context of memory, the hippocampus-inspired differentiable neural computer~\citep{graves2016hybrid} and its successors demonstrated how external memory with learned read/write operations can enhance neural network capabilities.

Epigenetics specifically has influenced AI research primarily in bioinformatics applications---predicting DNA methylation patterns~\citep{angermueller2017deepcpg}, modeling chromatin states~\citep{linder2025predicting}, and understanding gene regulatory networks~\citep{gao2024epigept}. However, the conceptual framework of epigenetics---dynamic regulation of fixed genetic information in response to environmental conditions---has not been systematically applied to AI system design.

\EpiContext{} represents the first application of epigenetic principles to agent context management. The mapping is both metaphorically rich and technically precise: the context graph corresponds to the genome, epigenetic tags to methylation/acetylation marks, and the fitness function to natural selection pressure. This cross-disciplinary bridge opens new avenues for both theoretical analysis and practical optimization of agent systems.

\subsection{Agent Reasoning and Planning}

Beyond memory, recent work has advanced agent reasoning capabilities through structured planning. Tree-of-Thoughts~\citep{yao2023tree} and Graph-of-Thoughts~\citep{besta2024graph} introduced structured reasoning topologies. ReCAP~\citep{recap2025neurips} proposed recursive context-aware reasoning with plan-ahead decomposition. MemAgent~\citep{yu2026memagent} used multi-conversation RL to train memory operations end-to-end. These approaches are complementary to \EpiContext: they focus on \textit{how} agents reason, while \EpiContext{} focuses on \textit{what information} agents reason with. The two directions can be combined synergistically.